
\documentclass{article}
\usepackage{colortbl}
\usepackage{makecell}
\usepackage{multirow}
\usepackage{supertabular}

\begin{document}

\newcounter{utterance}

\centering \large Interaction Transcript for game `codenames', experiment `risk\_low', episode 2 with mistral{-}large{-}3{-}azure.
\vspace{24pt}

{ \footnotesize  \setcounter{utterance}{1}
\setlength{\tabcolsep}{0pt}
\begin{supertabular}{c@{$\;$}|p{.15\linewidth}@{}p{.15\linewidth}p{.15\linewidth}p{.15\linewidth}p{.15\linewidth}p{.15\linewidth}}
    \# & $\;$A & \multicolumn{4}{c}{Game Master} & $\;\:$B\\
    \hline

    \theutterance \stepcounter{utterance}  
    & & & \multicolumn{2}{p{0.3\linewidth}}{
        \cellcolor[rgb]{0.9,0.9,0.9}{
            \makecell[{{p{\linewidth}}}]{
                \texttt{\tiny{[GM$|$GM]}}
                {'hidden': {'team': ['день', 'хмарочос', 'басейн', 'ягода', 'фільм', 'аварія', 'місяць', 'яблуко', 'машина'], 'innocent': ['виделка', 'вашингтон', 'кухоль', 'лимон', 'амазонка', 'стіл', 'пістолет', 'кенгуру'], 'opponent': ['душа', 'холод', 'батарея', 'горіх', 'скриня', 'африка', 'нота', 'мікроскоп'], 'assassin': []}, 'revealed': {'team': {'team': [], 'innocent': [], 'opponent': [], 'assassin': []}, 'opponent': {'team': [], 'innocent': [], 'opponent': [], 'assassin': []}}}
            }
        }
    }
    & & \\ \\

    \theutterance \stepcounter{utterance}  
    & & \multicolumn{4}{p{0.6\linewidth}}{
        \cellcolor[rgb]{0.9,0.9,0.9}{
            \makecell[{{p{\linewidth}}}]{
                \texttt{\tiny{[P1$\langle$GM]}}
                \texttt{Надайте одну однословну підказку, яка пов'язана з деякими словами вашої команди зі списку нижче. Ви можете обрати, зі скількома і з якими саме словами пов'язана ваша підказка.} \\
\texttt{Мета полягає в тому, щоб ваш партнер відгадав усі слова вашої команди на основі підказок, які ви йому даєте, раніше, ніж інша команда відкриє всі свої. Якщо ваш партнер відгадає слово «суперника» або «відволікача», його хід відгадування закінчується, і команда суперників отримує право зробити хід. Якщо ваш партнер відгадає слово «вбивці», гра закінчується, і ваша команда негайно програє.} \\
\texttt{Чим із більшою кількістю слів пов'язана ваша підказка, тим швидше ви можете виграти гру, але тим складнішими можуть бути асоціації для відгадування вашим партнером, тож обирайте слово{-}підказку мудро. Слово{-}підказка має бути семантично пов'язаним із цільовими словами; воно не може бути одним зі слів у списках або містити їхні частини.} \\
\\ 
\texttt{Завжди надавайте вашу однословну підказку та ваш розділений комами список пов'язаних цільових слів у такому форматі й робіть відповіді якомога коротшими; ніколи не додавайте жодного іншого тексту, окрім того, що потрібно в цій формі:} \\
\\ 
\texttt{ПІДКАЗКА: <СЛОВО>} \\
\texttt{ЦІЛІ: <ЦІЛІ>} \\
\\ 
\texttt{Слова вашої команди: день, хмарочос, басейн, ягода, фільм, аварія, місяць, яблуко, машина.} \\
\texttt{Слова команди суперника: душа, холод, батарея, горіх, скриня, африка, нота, мікроскоп.} \\
\texttt{Слова{-}відволікачі: виделка, вашингтон, кухоль, лимон, амазонка, стіл, пістолет, кенгуру.} \\
\texttt{Слова{-}вбивці: .} \\
            }
        }
    }
    & \\ \\

    \theutterance \stepcounter{utterance}  
    & \multicolumn{4}{p{0.6\linewidth}}{
        \cellcolor[rgb]{0.9,0.9,0.9}{
            \makecell[{{p{\linewidth}}}]{
                \texttt{\tiny{[P1$\rangle$GM]}}
                \texttt{ПІДКАЗКА: ніч} \\
\texttt{ЦІЛІ: день, місяць, фільм} \\
            }
        }
    }
    & & \\ \\

    \theutterance \stepcounter{utterance}  
    & & & \multicolumn{2}{p{0.3\linewidth}}{
        \cellcolor[rgb]{0.9,0.9,0.9}{
            \makecell[{{p{\linewidth}}}]{
                \texttt{\tiny{[GM$|$GM]}}
                \texttt{ніч} \\
            }
        }
    }
    & & \\ \\

    \theutterance \stepcounter{utterance}  
    & & & \multicolumn{2}{p{0.3\linewidth}}{
        \cellcolor[rgb]{0.9,0.9,0.9}{
            \makecell[{{p{\linewidth}}}]{
                \texttt{\tiny{[GM$|$GM]}}
                ['день', 'місяць', 'фільм']
            }
        }
    }
    & & \\ \\

    \theutterance \stepcounter{utterance}  
    & & & \multicolumn{2}{p{0.3\linewidth}}{
        \cellcolor[rgb]{0.9,0.9,0.9}{
            \makecell[{{p{\linewidth}}}]{
                \texttt{\tiny{[GM$|$GM]}}
                {'word': 'день', 'assignment': 'team'}
            }
        }
    }
    & & \\ \\

    \theutterance \stepcounter{utterance}  
    & & & \multicolumn{2}{p{0.3\linewidth}}{
        \cellcolor[rgb]{0.9,0.9,0.9}{
            \makecell[{{p{\linewidth}}}]{
                \texttt{\tiny{[GM$|$GM]}}
                {'word': 'місяць', 'assignment': 'team'}
            }
        }
    }
    & & \\ \\

    \theutterance \stepcounter{utterance}  
    & & & \multicolumn{2}{p{0.3\linewidth}}{
        \cellcolor[rgb]{0.9,0.9,0.9}{
            \makecell[{{p{\linewidth}}}]{
                \texttt{\tiny{[GM$|$GM]}}
                {'word': 'фільм', 'assignment': 'team'}
            }
        }
    }
    & & \\ \\

    \theutterance \stepcounter{utterance}  
    & & \multicolumn{4}{p{0.6\linewidth}}{
        \cellcolor[rgb]{0.9,0.9,0.9}{
            \makecell[{{p{\linewidth}}}]{
                \texttt{\tiny{[GM$\rangle$P2]}}
                \texttt{Надайте розділений комами список із не більш ніж 3 слів із наведеного нижче списку, які найкраще пов'язані або найбільш тісно асоціюються зі словом «ніч». Завжди починайте свій список здогадок із ВІДГАДКА:  і не додавайте жодного іншого тексту у своїй відповіді.} \\
\\ 
\texttt{аварія, холод, басейн, кухоль, лимон, вашингтон, горіх, батарея, амазонка, скриня, нота, день, яблуко, фільм, кенгуру, стіл, виделка, місяць, душа, мікроскоп, африка, машина, хмарочос, ягода, пістолет} \\
            }
        }
    }
    & \\ \\

    \theutterance \stepcounter{utterance}  
    & & & \multicolumn{4}{p{0.6\linewidth}}{
        \cellcolor[rgb]{0.9,0.9,0.9}{
            \makecell[{{p{\linewidth}}}]{
                \texttt{\tiny{[GM$\langle$P2]}}
                \texttt{ВІДГАДКА: місяць, холод, душа} \\
            }
        }
    }
     \\ \\

    \theutterance \stepcounter{utterance}  
    & & & \multicolumn{2}{p{0.3\linewidth}}{
        \cellcolor[rgb]{0.9,0.9,0.9}{
            \makecell[{{p{\linewidth}}}]{
                \texttt{\tiny{[GM$|$GM]}}
                ['місяць', 'холод', 'душа']
            }
        }
    }
    & & \\ \\

    \theutterance \stepcounter{utterance}  
    & & & \multicolumn{2}{p{0.3\linewidth}}{
        \cellcolor[rgb]{0.9,0.9,0.9}{
            \makecell[{{p{\linewidth}}}]{
                \texttt{\tiny{[GM$|$GM]}}
                {'word': 'місяць', 'assignment': 'team'}
            }
        }
    }
    & & \\ \\

    \theutterance \stepcounter{utterance}  
    & & & \multicolumn{2}{p{0.3\linewidth}}{
        \cellcolor[rgb]{0.9,0.9,0.9}{
            \makecell[{{p{\linewidth}}}]{
                \texttt{\tiny{[GM$|$GM]}}
                {'word': 'місяць', 'assignment': 'team'}
            }
        }
    }
    & & \\ \\

    \theutterance \stepcounter{utterance}  
    & & & \multicolumn{2}{p{0.3\linewidth}}{
        \cellcolor[rgb]{0.9,0.9,0.9}{
            \makecell[{{p{\linewidth}}}]{
                \texttt{\tiny{[GM$|$GM]}}
                {'word': 'холод', 'assignment': 'opponent'}
            }
        }
    }
    & & \\ \\

    \theutterance \stepcounter{utterance}  
    & & & \multicolumn{2}{p{0.3\linewidth}}{
        \cellcolor[rgb]{0.9,0.9,0.9}{
            \makecell[{{p{\linewidth}}}]{
                \texttt{\tiny{[GM$|$GM]}}
                \texttt{холод} \\
            }
        }
    }
    & & \\ \\

    \theutterance \stepcounter{utterance}  
    & & & \multicolumn{2}{p{0.3\linewidth}}{
        \cellcolor[rgb]{0.9,0.9,0.9}{
            \makecell[{{p{\linewidth}}}]{
                \texttt{\tiny{[GM$|$GM]}}
                {'word': 'батарея', 'assignment': 'opponent'}
            }
        }
    }
    & & \\ \\

    \theutterance \stepcounter{utterance}  
    & & & \multicolumn{2}{p{0.3\linewidth}}{
        \cellcolor[rgb]{0.9,0.9,0.9}{
            \makecell[{{p{\linewidth}}}]{
                \texttt{\tiny{[GM$|$GM]}}
                {'hidden': {'team': ['день', 'хмарочос', 'басейн', 'ягода', 'фільм', 'аварія', 'яблуко', 'машина'], 'innocent': ['виделка', 'вашингтон', 'кухоль', 'лимон', 'амазонка', 'стіл', 'пістолет', 'кенгуру'], 'opponent': ['душа', 'горіх', 'скриня', 'африка', 'нота', 'мікроскоп'], 'assassin': []}, 'revealed': {'team': {'team': ['місяць'], 'innocent': [], 'opponent': ['холод'], 'assassin': []}, 'opponent': {'team': [], 'innocent': [], 'opponent': ['батарея'], 'assassin': []}}}
            }
        }
    }
    & & \\ \\

    \theutterance \stepcounter{utterance}  
    & & \multicolumn{4}{p{0.6\linewidth}}{
        \cellcolor[rgb]{0.9,0.9,0.9}{
            \makecell[{{p{\linewidth}}}]{
                \texttt{\tiny{[P1$\langle$GM]}}
                \texttt{Слово місяць було відгадано правильно. Слово холод було відгадано, але це суперницьке слово. Хід вашого товариша по команді закінчився.} \\
\texttt{Тепер надайте ще одну підказку, пов'язану з деякими з ваших командних слів, що залишилися, та список пов'язаних цільових слів. Пам'ятайте: почніть підказку з 'ПІДКАЗКА: ', перейдіть на новий рядок і почніть ваш список цільових слів, розділених комами, з 'ЦІЛІ: '. Зверніть увагу: деякі слова було вилучено зі списків порівняно з попередніми запитами.} \\
\\ 
\texttt{Ваші командні слова, що залишилися: день, хмарочос, басейн, ягода, фільм, аварія, яблуко, машина.} \\
\texttt{Слова, що залишилися для вашого суперника: душа, батарея, горіх, скриня, африка, нота, мікроскоп.} \\
\texttt{Слова{-}відволікачі, що залишилися: виделка, вашингтон, кухоль, лимон, амазонка, стіл, пістолет, кенгуру.} \\
\texttt{Слова{-}вбивці, що залишилися: .} \\
            }
        }
    }
    & \\ \\

    \theutterance \stepcounter{utterance}  
    & \multicolumn{4}{p{0.6\linewidth}}{
        \cellcolor[rgb]{0.9,0.9,0.9}{
            \makecell[{{p{\linewidth}}}]{
                \texttt{\tiny{[P1$\rangle$GM]}}
                \texttt{ПІДКАЗКА: місто} \\
\texttt{ЦІЛІ: хмарочос, машина, аварія} \\
            }
        }
    }
    & & \\ \\

    \theutterance \stepcounter{utterance}  
    & & & \multicolumn{2}{p{0.3\linewidth}}{
        \cellcolor[rgb]{0.9,0.9,0.9}{
            \makecell[{{p{\linewidth}}}]{
                \texttt{\tiny{[GM$|$GM]}}
                \texttt{місто} \\
            }
        }
    }
    & & \\ \\

    \theutterance \stepcounter{utterance}  
    & & & \multicolumn{2}{p{0.3\linewidth}}{
        \cellcolor[rgb]{0.9,0.9,0.9}{
            \makecell[{{p{\linewidth}}}]{
                \texttt{\tiny{[GM$|$GM]}}
                ['хмарочос', 'машина', 'аварія']
            }
        }
    }
    & & \\ \\

    \theutterance \stepcounter{utterance}  
    & & & \multicolumn{2}{p{0.3\linewidth}}{
        \cellcolor[rgb]{0.9,0.9,0.9}{
            \makecell[{{p{\linewidth}}}]{
                \texttt{\tiny{[GM$|$GM]}}
                {'word': 'хмарочос', 'assignment': 'team'}
            }
        }
    }
    & & \\ \\

    \theutterance \stepcounter{utterance}  
    & & & \multicolumn{2}{p{0.3\linewidth}}{
        \cellcolor[rgb]{0.9,0.9,0.9}{
            \makecell[{{p{\linewidth}}}]{
                \texttt{\tiny{[GM$|$GM]}}
                {'word': 'машина', 'assignment': 'team'}
            }
        }
    }
    & & \\ \\

    \theutterance \stepcounter{utterance}  
    & & & \multicolumn{2}{p{0.3\linewidth}}{
        \cellcolor[rgb]{0.9,0.9,0.9}{
            \makecell[{{p{\linewidth}}}]{
                \texttt{\tiny{[GM$|$GM]}}
                {'word': 'аварія', 'assignment': 'team'}
            }
        }
    }
    & & \\ \\

    \theutterance \stepcounter{utterance}  
    & & \multicolumn{4}{p{0.6\linewidth}}{
        \cellcolor[rgb]{0.9,0.9,0.9}{
            \makecell[{{p{\linewidth}}}]{
                \texttt{\tiny{[GM$\rangle$P2]}}
                \texttt{Слово місяць було відгадано правильно. Слово холод було відгадано, але це суперницьке слово. Ваш хід закінчився.} \\
\texttt{Тепер надайте ще один список слів, розділених комами, щонайменше з 1 і щонайбільше з 3 слів із наведеного нижче списку слів, які найкраще пов'язані або найбільш тісно асоціюються зі словом 'місто'. Пам'ятайте, що потрібно почати вашу відповідь з ВІДГАДКА: . Зауваження: деякі слова було вилучено зі списку порівняно з попередніми запитами.} \\
\\ 
\texttt{аварія, басейн, кухоль, лимон, вашингтон, горіх, амазонка, скриня, нота, день, яблуко, фільм, кенгуру, стіл, виделка, душа, мікроскоп, африка, машина, хмарочос, ягода, пістолет} \\
            }
        }
    }
    & \\ \\

    \theutterance \stepcounter{utterance}  
    & & & \multicolumn{4}{p{0.6\linewidth}}{
        \cellcolor[rgb]{0.9,0.9,0.9}{
            \makecell[{{p{\linewidth}}}]{
                \texttt{\tiny{[GM$\langle$P2]}}
                \texttt{ВІДГАДКА: хмарочос, машина, вашингтон} \\
            }
        }
    }
     \\ \\

    \theutterance \stepcounter{utterance}  
    & & & \multicolumn{2}{p{0.3\linewidth}}{
        \cellcolor[rgb]{0.9,0.9,0.9}{
            \makecell[{{p{\linewidth}}}]{
                \texttt{\tiny{[GM$|$GM]}}
                ['хмарочос', 'машина', 'вашингтон']
            }
        }
    }
    & & \\ \\

    \theutterance \stepcounter{utterance}  
    & & & \multicolumn{2}{p{0.3\linewidth}}{
        \cellcolor[rgb]{0.9,0.9,0.9}{
            \makecell[{{p{\linewidth}}}]{
                \texttt{\tiny{[GM$|$GM]}}
                {'word': 'хмарочос', 'assignment': 'team'}
            }
        }
    }
    & & \\ \\

    \theutterance \stepcounter{utterance}  
    & & & \multicolumn{2}{p{0.3\linewidth}}{
        \cellcolor[rgb]{0.9,0.9,0.9}{
            \makecell[{{p{\linewidth}}}]{
                \texttt{\tiny{[GM$|$GM]}}
                {'word': 'хмарочос', 'assignment': 'team'}
            }
        }
    }
    & & \\ \\

    \theutterance \stepcounter{utterance}  
    & & & \multicolumn{2}{p{0.3\linewidth}}{
        \cellcolor[rgb]{0.9,0.9,0.9}{
            \makecell[{{p{\linewidth}}}]{
                \texttt{\tiny{[GM$|$GM]}}
                {'word': 'машина', 'assignment': 'team'}
            }
        }
    }
    & & \\ \\

    \theutterance \stepcounter{utterance}  
    & & & \multicolumn{2}{p{0.3\linewidth}}{
        \cellcolor[rgb]{0.9,0.9,0.9}{
            \makecell[{{p{\linewidth}}}]{
                \texttt{\tiny{[GM$|$GM]}}
                {'word': 'машина', 'assignment': 'team'}
            }
        }
    }
    & & \\ \\

    \theutterance \stepcounter{utterance}  
    & & & \multicolumn{2}{p{0.3\linewidth}}{
        \cellcolor[rgb]{0.9,0.9,0.9}{
            \makecell[{{p{\linewidth}}}]{
                \texttt{\tiny{[GM$|$GM]}}
                {'word': 'вашингтон', 'assignment': 'innocent'}
            }
        }
    }
    & & \\ \\

    \theutterance \stepcounter{utterance}  
    & & & \multicolumn{2}{p{0.3\linewidth}}{
        \cellcolor[rgb]{0.9,0.9,0.9}{
            \makecell[{{p{\linewidth}}}]{
                \texttt{\tiny{[GM$|$GM]}}
                \texttt{вашингтон} \\
            }
        }
    }
    & & \\ \\

    \theutterance \stepcounter{utterance}  
    & & & \multicolumn{2}{p{0.3\linewidth}}{
        \cellcolor[rgb]{0.9,0.9,0.9}{
            \makecell[{{p{\linewidth}}}]{
                \texttt{\tiny{[GM$|$GM]}}
                {'word': 'горіх', 'assignment': 'opponent'}
            }
        }
    }
    & & \\ \\

    \theutterance \stepcounter{utterance}  
    & & & \multicolumn{2}{p{0.3\linewidth}}{
        \cellcolor[rgb]{0.9,0.9,0.9}{
            \makecell[{{p{\linewidth}}}]{
                \texttt{\tiny{[GM$|$GM]}}
                {'hidden': {'team': ['день', 'басейн', 'ягода', 'фільм', 'аварія', 'яблуко'], 'innocent': ['виделка', 'кухоль', 'лимон', 'амазонка', 'стіл', 'пістолет', 'кенгуру'], 'opponent': ['душа', 'скриня', 'африка', 'нота', 'мікроскоп'], 'assassin': []}, 'revealed': {'team': {'team': ['місяць', 'хмарочос', 'машина'], 'innocent': ['вашингтон'], 'opponent': ['холод'], 'assassin': []}, 'opponent': {'team': [], 'innocent': [], 'opponent': ['батарея', 'горіх'], 'assassin': []}}}
            }
        }
    }
    & & \\ \\

    \theutterance \stepcounter{utterance}  
    & & \multicolumn{4}{p{0.6\linewidth}}{
        \cellcolor[rgb]{0.9,0.9,0.9}{
            \makecell[{{p{\linewidth}}}]{
                \texttt{\tiny{[P1$\langle$GM]}}
                \texttt{Слова хмарочос, машина було відгадано правильно. Слово вашингтон було відгадано, але це нейтральне слово. Хід вашого товариша по команді закінчився.} \\
\texttt{Тепер надайте ще одну підказку, пов'язану з деякими з ваших командних слів, що залишилися, та список пов'язаних цільових слів. Пам'ятайте: почніть підказку з 'ПІДКАЗКА: ', перейдіть на новий рядок і почніть ваш список цільових слів, розділених комами, з 'ЦІЛІ: '. Зверніть увагу: деякі слова було вилучено зі списків порівняно з попередніми запитами.} \\
\\ 
\texttt{Ваші командні слова, що залишилися: день, басейн, ягода, фільм, аварія, яблуко.} \\
\texttt{Слова, що залишилися для вашого суперника: душа, горіх, скриня, африка, нота, мікроскоп.} \\
\texttt{Слова{-}відволікачі, що залишилися: виделка, кухоль, лимон, амазонка, стіл, пістолет, кенгуру.} \\
\texttt{Слова{-}вбивці, що залишилися: .} \\
            }
        }
    }
    & \\ \\

    \theutterance \stepcounter{utterance}  
    & \multicolumn{4}{p{0.6\linewidth}}{
        \cellcolor[rgb]{0.9,0.9,0.9}{
            \makecell[{{p{\linewidth}}}]{
                \texttt{\tiny{[P1$\rangle$GM]}}
                \texttt{ПІДКАЗКА: літо} \\
\texttt{ЦІЛІ: день, басейн, ягода, яблуко} \\
            }
        }
    }
    & & \\ \\

    \theutterance \stepcounter{utterance}  
    & & & \multicolumn{2}{p{0.3\linewidth}}{
        \cellcolor[rgb]{0.9,0.9,0.9}{
            \makecell[{{p{\linewidth}}}]{
                \texttt{\tiny{[GM$|$GM]}}
                \texttt{літо} \\
            }
        }
    }
    & & \\ \\

    \theutterance \stepcounter{utterance}  
    & & & \multicolumn{2}{p{0.3\linewidth}}{
        \cellcolor[rgb]{0.9,0.9,0.9}{
            \makecell[{{p{\linewidth}}}]{
                \texttt{\tiny{[GM$|$GM]}}
                ['день', 'басейн', 'ягода', 'яблуко']
            }
        }
    }
    & & \\ \\

    \theutterance \stepcounter{utterance}  
    & & & \multicolumn{2}{p{0.3\linewidth}}{
        \cellcolor[rgb]{0.9,0.9,0.9}{
            \makecell[{{p{\linewidth}}}]{
                \texttt{\tiny{[GM$|$GM]}}
                {'word': 'день', 'assignment': 'team'}
            }
        }
    }
    & & \\ \\

    \theutterance \stepcounter{utterance}  
    & & & \multicolumn{2}{p{0.3\linewidth}}{
        \cellcolor[rgb]{0.9,0.9,0.9}{
            \makecell[{{p{\linewidth}}}]{
                \texttt{\tiny{[GM$|$GM]}}
                {'word': 'басейн', 'assignment': 'team'}
            }
        }
    }
    & & \\ \\

    \theutterance \stepcounter{utterance}  
    & & & \multicolumn{2}{p{0.3\linewidth}}{
        \cellcolor[rgb]{0.9,0.9,0.9}{
            \makecell[{{p{\linewidth}}}]{
                \texttt{\tiny{[GM$|$GM]}}
                {'word': 'ягода', 'assignment': 'team'}
            }
        }
    }
    & & \\ \\

    \theutterance \stepcounter{utterance}  
    & & & \multicolumn{2}{p{0.3\linewidth}}{
        \cellcolor[rgb]{0.9,0.9,0.9}{
            \makecell[{{p{\linewidth}}}]{
                \texttt{\tiny{[GM$|$GM]}}
                {'word': 'яблуко', 'assignment': 'team'}
            }
        }
    }
    & & \\ \\

    \theutterance \stepcounter{utterance}  
    & & \multicolumn{4}{p{0.6\linewidth}}{
        \cellcolor[rgb]{0.9,0.9,0.9}{
            \makecell[{{p{\linewidth}}}]{
                \texttt{\tiny{[GM$\rangle$P2]}}
                \texttt{Слова хмарочос, машина було відгадано правильно. Слово вашингтон було відгадано, але це нейтральне слово. Ваш хід закінчився.} \\
\texttt{Тепер надайте ще один список слів, розділених комами, щонайменше з 1 і щонайбільше з 4 слів із наведеного нижче списку слів, які найкраще пов'язані або найбільш тісно асоціюються зі словом 'літо'. Пам'ятайте, що потрібно почати вашу відповідь з ВІДГАДКА: . Зауваження: деякі слова було вилучено зі списку порівняно з попередніми запитами.} \\
\\ 
\texttt{аварія, басейн, кухоль, лимон, амазонка, скриня, нота, день, яблуко, фільм, кенгуру, стіл, виделка, душа, мікроскоп, африка, ягода, пістолет} \\
            }
        }
    }
    & \\ \\

    \theutterance \stepcounter{utterance}  
    & & & \multicolumn{4}{p{0.6\linewidth}}{
        \cellcolor[rgb]{0.9,0.9,0.9}{
            \makecell[{{p{\linewidth}}}]{
                \texttt{\tiny{[GM$\langle$P2]}}
                \texttt{ВІДГАДКА: басейн, ягода, день, лимон} \\
            }
        }
    }
     \\ \\

    \theutterance \stepcounter{utterance}  
    & & & \multicolumn{2}{p{0.3\linewidth}}{
        \cellcolor[rgb]{0.9,0.9,0.9}{
            \makecell[{{p{\linewidth}}}]{
                \texttt{\tiny{[GM$|$GM]}}
                ['басейн', 'ягода', 'день', 'лимон']
            }
        }
    }
    & & \\ \\

    \theutterance \stepcounter{utterance}  
    & & & \multicolumn{2}{p{0.3\linewidth}}{
        \cellcolor[rgb]{0.9,0.9,0.9}{
            \makecell[{{p{\linewidth}}}]{
                \texttt{\tiny{[GM$|$GM]}}
                {'word': 'басейн', 'assignment': 'team'}
            }
        }
    }
    & & \\ \\

    \theutterance \stepcounter{utterance}  
    & & & \multicolumn{2}{p{0.3\linewidth}}{
        \cellcolor[rgb]{0.9,0.9,0.9}{
            \makecell[{{p{\linewidth}}}]{
                \texttt{\tiny{[GM$|$GM]}}
                {'word': 'басейн', 'assignment': 'team'}
            }
        }
    }
    & & \\ \\

    \theutterance \stepcounter{utterance}  
    & & & \multicolumn{2}{p{0.3\linewidth}}{
        \cellcolor[rgb]{0.9,0.9,0.9}{
            \makecell[{{p{\linewidth}}}]{
                \texttt{\tiny{[GM$|$GM]}}
                {'word': 'ягода', 'assignment': 'team'}
            }
        }
    }
    & & \\ \\

    \theutterance \stepcounter{utterance}  
    & & & \multicolumn{2}{p{0.3\linewidth}}{
        \cellcolor[rgb]{0.9,0.9,0.9}{
            \makecell[{{p{\linewidth}}}]{
                \texttt{\tiny{[GM$|$GM]}}
                {'word': 'ягода', 'assignment': 'team'}
            }
        }
    }
    & & \\ \\

    \theutterance \stepcounter{utterance}  
    & & & \multicolumn{2}{p{0.3\linewidth}}{
        \cellcolor[rgb]{0.9,0.9,0.9}{
            \makecell[{{p{\linewidth}}}]{
                \texttt{\tiny{[GM$|$GM]}}
                {'word': 'день', 'assignment': 'team'}
            }
        }
    }
    & & \\ \\

    \theutterance \stepcounter{utterance}  
    & & & \multicolumn{2}{p{0.3\linewidth}}{
        \cellcolor[rgb]{0.9,0.9,0.9}{
            \makecell[{{p{\linewidth}}}]{
                \texttt{\tiny{[GM$|$GM]}}
                {'word': 'день', 'assignment': 'team'}
            }
        }
    }
    & & \\ \\

    \theutterance \stepcounter{utterance}  
    & & & \multicolumn{2}{p{0.3\linewidth}}{
        \cellcolor[rgb]{0.9,0.9,0.9}{
            \makecell[{{p{\linewidth}}}]{
                \texttt{\tiny{[GM$|$GM]}}
                {'word': 'лимон', 'assignment': 'innocent'}
            }
        }
    }
    & & \\ \\

    \theutterance \stepcounter{utterance}  
    & & & \multicolumn{2}{p{0.3\linewidth}}{
        \cellcolor[rgb]{0.9,0.9,0.9}{
            \makecell[{{p{\linewidth}}}]{
                \texttt{\tiny{[GM$|$GM]}}
                \texttt{лимон} \\
            }
        }
    }
    & & \\ \\

    \theutterance \stepcounter{utterance}  
    & & & \multicolumn{2}{p{0.3\linewidth}}{
        \cellcolor[rgb]{0.9,0.9,0.9}{
            \makecell[{{p{\linewidth}}}]{
                \texttt{\tiny{[GM$|$GM]}}
                {'word': 'душа', 'assignment': 'opponent'}
            }
        }
    }
    & & \\ \\

    \theutterance \stepcounter{utterance}  
    & & & \multicolumn{2}{p{0.3\linewidth}}{
        \cellcolor[rgb]{0.9,0.9,0.9}{
            \makecell[{{p{\linewidth}}}]{
                \texttt{\tiny{[GM$|$GM]}}
                {'hidden': {'team': ['фільм', 'аварія', 'яблуко'], 'innocent': ['виделка', 'кухоль', 'амазонка', 'стіл', 'пістолет', 'кенгуру'], 'opponent': ['скриня', 'африка', 'нота', 'мікроскоп'], 'assassin': []}, 'revealed': {'team': {'team': ['місяць', 'хмарочос', 'машина', 'басейн', 'ягода', 'день'], 'innocent': ['вашингтон', 'лимон'], 'opponent': ['холод'], 'assassin': []}, 'opponent': {'team': [], 'innocent': [], 'opponent': ['батарея', 'горіх', 'душа'], 'assassin': []}}}
            }
        }
    }
    & & \\ \\

    \theutterance \stepcounter{utterance}  
    & & \multicolumn{4}{p{0.6\linewidth}}{
        \cellcolor[rgb]{0.9,0.9,0.9}{
            \makecell[{{p{\linewidth}}}]{
                \texttt{\tiny{[P1$\langle$GM]}}
                \texttt{Слова басейн, ягода, день було відгадано правильно. Слово лимон було відгадано, але це нейтральне слово. Хід вашого товариша по команді закінчився.} \\
\texttt{Тепер надайте ще одну підказку, пов'язану з деякими з ваших командних слів, що залишилися, та список пов'язаних цільових слів. Пам'ятайте: почніть підказку з 'ПІДКАЗКА: ', перейдіть на новий рядок і почніть ваш список цільових слів, розділених комами, з 'ЦІЛІ: '. Зверніть увагу: деякі слова було вилучено зі списків порівняно з попередніми запитами.} \\
\\ 
\texttt{Ваші командні слова, що залишилися: фільм, аварія, яблуко.} \\
\texttt{Слова, що залишилися для вашого суперника: душа, скриня, африка, нота, мікроскоп.} \\
\texttt{Слова{-}відволікачі, що залишилися: виделка, кухоль, амазонка, стіл, пістолет, кенгуру.} \\
\texttt{Слова{-}вбивці, що залишилися: .} \\
            }
        }
    }
    & \\ \\

    \theutterance \stepcounter{utterance}  
    & \multicolumn{4}{p{0.6\linewidth}}{
        \cellcolor[rgb]{0.9,0.9,0.9}{
            \makecell[{{p{\linewidth}}}]{
                \texttt{\tiny{[P1$\rangle$GM]}}
                \texttt{ПІДКАЗКА: екран} \\
\texttt{ЦІЛІ: фільм, аварія} \\
            }
        }
    }
    & & \\ \\

    \theutterance \stepcounter{utterance}  
    & & & \multicolumn{2}{p{0.3\linewidth}}{
        \cellcolor[rgb]{0.9,0.9,0.9}{
            \makecell[{{p{\linewidth}}}]{
                \texttt{\tiny{[GM$|$GM]}}
                \texttt{екран} \\
            }
        }
    }
    & & \\ \\

    \theutterance \stepcounter{utterance}  
    & & & \multicolumn{2}{p{0.3\linewidth}}{
        \cellcolor[rgb]{0.9,0.9,0.9}{
            \makecell[{{p{\linewidth}}}]{
                \texttt{\tiny{[GM$|$GM]}}
                ['фільм', 'аварія']
            }
        }
    }
    & & \\ \\

    \theutterance \stepcounter{utterance}  
    & & & \multicolumn{2}{p{0.3\linewidth}}{
        \cellcolor[rgb]{0.9,0.9,0.9}{
            \makecell[{{p{\linewidth}}}]{
                \texttt{\tiny{[GM$|$GM]}}
                {'word': 'фільм', 'assignment': 'team'}
            }
        }
    }
    & & \\ \\

    \theutterance \stepcounter{utterance}  
    & & & \multicolumn{2}{p{0.3\linewidth}}{
        \cellcolor[rgb]{0.9,0.9,0.9}{
            \makecell[{{p{\linewidth}}}]{
                \texttt{\tiny{[GM$|$GM]}}
                {'word': 'аварія', 'assignment': 'team'}
            }
        }
    }
    & & \\ \\

    \theutterance \stepcounter{utterance}  
    & & \multicolumn{4}{p{0.6\linewidth}}{
        \cellcolor[rgb]{0.9,0.9,0.9}{
            \makecell[{{p{\linewidth}}}]{
                \texttt{\tiny{[GM$\rangle$P2]}}
                \texttt{Слова басейн, ягода, день було відгадано правильно. Слово лимон було відгадано, але це нейтральне слово. Ваш хід закінчився.} \\
\texttt{Тепер надайте ще один список слів, розділених комами, щонайменше з 1 і щонайбільше з 2 слів із наведеного нижче списку слів, які найкраще пов'язані або найбільш тісно асоціюються зі словом 'екран'. Пам'ятайте, що потрібно почати вашу відповідь з ВІДГАДКА: . Зауваження: деякі слова було вилучено зі списку порівняно з попередніми запитами.} \\
\\ 
\texttt{аварія, кухоль, амазонка, скриня, нота, яблуко, фільм, кенгуру, стіл, виделка, мікроскоп, африка, пістолет} \\
            }
        }
    }
    & \\ \\

    \theutterance \stepcounter{utterance}  
    & & & \multicolumn{4}{p{0.6\linewidth}}{
        \cellcolor[rgb]{0.9,0.9,0.9}{
            \makecell[{{p{\linewidth}}}]{
                \texttt{\tiny{[GM$\langle$P2]}}
                \texttt{ВІДГАДКА: фільм, мікроскоп} \\
            }
        }
    }
     \\ \\

    \theutterance \stepcounter{utterance}  
    & & & \multicolumn{2}{p{0.3\linewidth}}{
        \cellcolor[rgb]{0.9,0.9,0.9}{
            \makecell[{{p{\linewidth}}}]{
                \texttt{\tiny{[GM$|$GM]}}
                ['фільм', 'мікроскоп']
            }
        }
    }
    & & \\ \\

    \theutterance \stepcounter{utterance}  
    & & & \multicolumn{2}{p{0.3\linewidth}}{
        \cellcolor[rgb]{0.9,0.9,0.9}{
            \makecell[{{p{\linewidth}}}]{
                \texttt{\tiny{[GM$|$GM]}}
                {'word': 'фільм', 'assignment': 'team'}
            }
        }
    }
    & & \\ \\

    \theutterance \stepcounter{utterance}  
    & & & \multicolumn{2}{p{0.3\linewidth}}{
        \cellcolor[rgb]{0.9,0.9,0.9}{
            \makecell[{{p{\linewidth}}}]{
                \texttt{\tiny{[GM$|$GM]}}
                {'word': 'фільм', 'assignment': 'team'}
            }
        }
    }
    & & \\ \\

    \theutterance \stepcounter{utterance}  
    & & & \multicolumn{2}{p{0.3\linewidth}}{
        \cellcolor[rgb]{0.9,0.9,0.9}{
            \makecell[{{p{\linewidth}}}]{
                \texttt{\tiny{[GM$|$GM]}}
                {'word': 'мікроскоп', 'assignment': 'opponent'}
            }
        }
    }
    & & \\ \\

    \theutterance \stepcounter{utterance}  
    & & & \multicolumn{2}{p{0.3\linewidth}}{
        \cellcolor[rgb]{0.9,0.9,0.9}{
            \makecell[{{p{\linewidth}}}]{
                \texttt{\tiny{[GM$|$GM]}}
                \texttt{мікроскоп} \\
            }
        }
    }
    & & \\ \\

    \theutterance \stepcounter{utterance}  
    & & & \multicolumn{2}{p{0.3\linewidth}}{
        \cellcolor[rgb]{0.9,0.9,0.9}{
            \makecell[{{p{\linewidth}}}]{
                \texttt{\tiny{[GM$|$GM]}}
                {'word': 'нота', 'assignment': 'opponent'}
            }
        }
    }
    & & \\ \\

    \theutterance \stepcounter{utterance}  
    & & & \multicolumn{2}{p{0.3\linewidth}}{
        \cellcolor[rgb]{0.9,0.9,0.9}{
            \makecell[{{p{\linewidth}}}]{
                \texttt{\tiny{[GM$|$GM]}}
                {'hidden': {'team': ['аварія', 'яблуко'], 'innocent': ['виделка', 'кухоль', 'амазонка', 'стіл', 'пістолет', 'кенгуру'], 'opponent': ['скриня', 'африка'], 'assassin': []}, 'revealed': {'team': {'team': ['місяць', 'хмарочос', 'машина', 'басейн', 'ягода', 'день', 'фільм'], 'innocent': ['вашингтон', 'лимон'], 'opponent': ['холод', 'мікроскоп'], 'assassin': []}, 'opponent': {'team': [], 'innocent': [], 'opponent': ['батарея', 'горіх', 'душа', 'нота'], 'assassin': []}}}
            }
        }
    }
    & & \\ \\

    \theutterance \stepcounter{utterance}  
    & & \multicolumn{4}{p{0.6\linewidth}}{
        \cellcolor[rgb]{0.9,0.9,0.9}{
            \makecell[{{p{\linewidth}}}]{
                \texttt{\tiny{[P1$\langle$GM]}}
                \texttt{Слово фільм було відгадано правильно. Слово мікроскоп було відгадано, але це суперницьке слово. Хід вашого товариша по команді закінчився.} \\
\texttt{Тепер надайте ще одну підказку, пов'язану з деякими з ваших командних слів, що залишилися, та список пов'язаних цільових слів. Пам'ятайте: почніть підказку з 'ПІДКАЗКА: ', перейдіть на новий рядок і почніть ваш список цільових слів, розділених комами, з 'ЦІЛІ: '. Зверніть увагу: деякі слова було вилучено зі списків порівняно з попередніми запитами.} \\
\\ 
\texttt{Ваші командні слова, що залишилися: аварія, яблуко.} \\
\texttt{Слова, що залишилися для вашого суперника: скриня, африка, нота.} \\
\texttt{Слова{-}відволікачі, що залишилися: виделка, кухоль, амазонка, стіл, пістолет, кенгуру.} \\
\texttt{Слова{-}вбивці, що залишилися: .} \\
            }
        }
    }
    & \\ \\

    \theutterance \stepcounter{utterance}  
    & \multicolumn{4}{p{0.6\linewidth}}{
        \cellcolor[rgb]{0.9,0.9,0.9}{
            \makecell[{{p{\linewidth}}}]{
                \texttt{\tiny{[P1$\rangle$GM]}}
                \texttt{ПІДКАЗКА: Ньютон} \\
\texttt{ЦІЛІ: яблуко, аварія} \\
            }
        }
    }
    & & \\ \\

    \theutterance \stepcounter{utterance}  
    & & & \multicolumn{2}{p{0.3\linewidth}}{
        \cellcolor[rgb]{0.9,0.9,0.9}{
            \makecell[{{p{\linewidth}}}]{
                \texttt{\tiny{[GM$|$GM]}}
                \texttt{ньютон} \\
            }
        }
    }
    & & \\ \\

    \theutterance \stepcounter{utterance}  
    & & & \multicolumn{2}{p{0.3\linewidth}}{
        \cellcolor[rgb]{0.9,0.9,0.9}{
            \makecell[{{p{\linewidth}}}]{
                \texttt{\tiny{[GM$|$GM]}}
                ['яблуко', 'аварія']
            }
        }
    }
    & & \\ \\

    \theutterance \stepcounter{utterance}  
    & & & \multicolumn{2}{p{0.3\linewidth}}{
        \cellcolor[rgb]{0.9,0.9,0.9}{
            \makecell[{{p{\linewidth}}}]{
                \texttt{\tiny{[GM$|$GM]}}
                {'word': 'яблуко', 'assignment': 'team'}
            }
        }
    }
    & & \\ \\

    \theutterance \stepcounter{utterance}  
    & & & \multicolumn{2}{p{0.3\linewidth}}{
        \cellcolor[rgb]{0.9,0.9,0.9}{
            \makecell[{{p{\linewidth}}}]{
                \texttt{\tiny{[GM$|$GM]}}
                {'word': 'аварія', 'assignment': 'team'}
            }
        }
    }
    & & \\ \\

    \theutterance \stepcounter{utterance}  
    & & \multicolumn{4}{p{0.6\linewidth}}{
        \cellcolor[rgb]{0.9,0.9,0.9}{
            \makecell[{{p{\linewidth}}}]{
                \texttt{\tiny{[GM$\rangle$P2]}}
                \texttt{Слово фільм було відгадано правильно. Слово мікроскоп було відгадано, але це суперницьке слово. Ваш хід закінчився.} \\
\texttt{Тепер надайте ще один список слів, розділених комами, щонайменше з 1 і щонайбільше з 2 слів із наведеного нижче списку слів, які найкраще пов'язані або найбільш тісно асоціюються зі словом 'ньютон'. Пам'ятайте, що потрібно почати вашу відповідь з ВІДГАДКА: . Зауваження: деякі слова було вилучено зі списку порівняно з попередніми запитами.} \\
\\ 
\texttt{аварія, кухоль, амазонка, скриня, яблуко, кенгуру, стіл, виделка, африка, пістолет} \\
            }
        }
    }
    & \\ \\

    \theutterance \stepcounter{utterance}  
    & & & \multicolumn{4}{p{0.6\linewidth}}{
        \cellcolor[rgb]{0.9,0.9,0.9}{
            \makecell[{{p{\linewidth}}}]{
                \texttt{\tiny{[GM$\langle$P2]}}
                \texttt{ВІДГАДКА: яблуко, скриня} \\
            }
        }
    }
     \\ \\

    \theutterance \stepcounter{utterance}  
    & & & \multicolumn{2}{p{0.3\linewidth}}{
        \cellcolor[rgb]{0.9,0.9,0.9}{
            \makecell[{{p{\linewidth}}}]{
                \texttt{\tiny{[GM$|$GM]}}
                ['яблуко', 'скриня']
            }
        }
    }
    & & \\ \\

    \theutterance \stepcounter{utterance}  
    & & & \multicolumn{2}{p{0.3\linewidth}}{
        \cellcolor[rgb]{0.9,0.9,0.9}{
            \makecell[{{p{\linewidth}}}]{
                \texttt{\tiny{[GM$|$GM]}}
                {'word': 'яблуко', 'assignment': 'team'}
            }
        }
    }
    & & \\ \\

    \theutterance \stepcounter{utterance}  
    & & & \multicolumn{2}{p{0.3\linewidth}}{
        \cellcolor[rgb]{0.9,0.9,0.9}{
            \makecell[{{p{\linewidth}}}]{
                \texttt{\tiny{[GM$|$GM]}}
                {'word': 'яблуко', 'assignment': 'team'}
            }
        }
    }
    & & \\ \\

    \theutterance \stepcounter{utterance}  
    & & & \multicolumn{2}{p{0.3\linewidth}}{
        \cellcolor[rgb]{0.9,0.9,0.9}{
            \makecell[{{p{\linewidth}}}]{
                \texttt{\tiny{[GM$|$GM]}}
                {'word': 'скриня', 'assignment': 'opponent'}
            }
        }
    }
    & & \\ \\

    \theutterance \stepcounter{utterance}  
    & & & \multicolumn{2}{p{0.3\linewidth}}{
        \cellcolor[rgb]{0.9,0.9,0.9}{
            \makecell[{{p{\linewidth}}}]{
                \texttt{\tiny{[GM$|$GM]}}
                \texttt{скриня} \\
            }
        }
    }
    & & \\ \\

    \theutterance \stepcounter{utterance}  
    & & & \multicolumn{2}{p{0.3\linewidth}}{
        \cellcolor[rgb]{0.9,0.9,0.9}{
            \makecell[{{p{\linewidth}}}]{
                \texttt{\tiny{[GM$|$GM]}}
                {'word': 'африка', 'assignment': 'opponent'}
            }
        }
    }
    & & \\ \\

    \theutterance \stepcounter{utterance}  
    & & & \multicolumn{2}{p{0.3\linewidth}}{
        \cellcolor[rgb]{0.9,0.9,0.9}{
            \makecell[{{p{\linewidth}}}]{
                \texttt{\tiny{[GM$|$GM]}}
                {'hidden': {'team': ['аварія'], 'innocent': ['виделка', 'кухоль', 'амазонка', 'стіл', 'пістолет', 'кенгуру'], 'opponent': [], 'assassin': []}, 'revealed': {'team': {'team': ['місяць', 'хмарочос', 'машина', 'басейн', 'ягода', 'день', 'фільм', 'яблуко'], 'innocent': ['вашингтон', 'лимон'], 'opponent': ['холод', 'мікроскоп', 'скриня'], 'assassin': []}, 'opponent': {'team': [], 'innocent': [], 'opponent': ['батарея', 'горіх', 'душа', 'нота', 'африка'], 'assassin': []}}}
            }
        }
    }
    & & \\ \\

    \theutterance \stepcounter{utterance}  
    & & \multicolumn{4}{p{0.6\linewidth}}{
        \cellcolor[rgb]{0.9,0.9,0.9}{
            \makecell[{{p{\linewidth}}}]{
                \texttt{\tiny{[P1$\langle$GM]}}
                \texttt{Слово яблуко було відгадано правильно. Слово скриня було відгадано, але це суперницьке слово. Хід вашого товариша по команді закінчився.} \\
\texttt{Тепер надайте ще одну підказку, пов'язану з деякими з ваших командних слів, що залишилися, та список пов'язаних цільових слів. Пам'ятайте: почніть підказку з 'ПІДКАЗКА: ', перейдіть на новий рядок і почніть ваш список цільових слів, розділених комами, з 'ЦІЛІ: '. Зверніть увагу: деякі слова було вилучено зі списків порівняно з попередніми запитами.} \\
\\ 
\texttt{Ваші командні слова, що залишилися: аварія.} \\
\texttt{Слова, що залишилися для вашого суперника: африка.} \\
\texttt{Слова{-}відволікачі, що залишилися: виделка, кухоль, амазонка, стіл, пістолет, кенгуру.} \\
\texttt{Слова{-}вбивці, що залишилися: .} \\
            }
        }
    }
    & \\ \\

    \theutterance \stepcounter{utterance}  
    & \multicolumn{4}{p{0.6\linewidth}}{
        \cellcolor[rgb]{0.9,0.9,0.9}{
            \makecell[{{p{\linewidth}}}]{
                \texttt{\tiny{[P1$\rangle$GM]}}
                \texttt{ПІДКАЗКА: ДТП} \\
\texttt{ЦІЛІ: аварія} \\
            }
        }
    }
    & & \\ \\

    \theutterance \stepcounter{utterance}  
    & & & \multicolumn{2}{p{0.3\linewidth}}{
        \cellcolor[rgb]{0.9,0.9,0.9}{
            \makecell[{{p{\linewidth}}}]{
                \texttt{\tiny{[GM$|$GM]}}
                \texttt{дтп} \\
            }
        }
    }
    & & \\ \\

    \theutterance \stepcounter{utterance}  
    & & & \multicolumn{2}{p{0.3\linewidth}}{
        \cellcolor[rgb]{0.9,0.9,0.9}{
            \makecell[{{p{\linewidth}}}]{
                \texttt{\tiny{[GM$|$GM]}}
                ['аварія']
            }
        }
    }
    & & \\ \\

    \theutterance \stepcounter{utterance}  
    & & & \multicolumn{2}{p{0.3\linewidth}}{
        \cellcolor[rgb]{0.9,0.9,0.9}{
            \makecell[{{p{\linewidth}}}]{
                \texttt{\tiny{[GM$|$GM]}}
                {'word': 'аварія', 'assignment': 'team'}
            }
        }
    }
    & & \\ \\

    \theutterance \stepcounter{utterance}  
    & & & \multicolumn{2}{p{0.3\linewidth}}{
        \cellcolor[rgb]{0.9,0.9,0.9}{
            \makecell[{{p{\linewidth}}}]{
                \texttt{\tiny{[GM$|$GM]}}
                \texttt{opponent has won} \\
            }
        }
    }
    & & \\ \\

\end{supertabular}
}

\end{document}
